\section{Datasets Seleccionados}

\subsection{Dataset 1: Yeast Protein-Protein Interaction Network}

\begin{center}
\begin{tabular}{ll}
\toprule
\textbf{Característica} & \textbf{Valor} \\
\midrule
Tipo & No dirigido, no ponderado \\
Vértices & 6.526 proteínas \\
Aristas & 532.180 interacciones \\
Formato & Lista de aristas separada por tabuladores \\
Fuente & Kaggle / Alexander Chervov \\
\bottomrule
\end{tabular}
\end{center}

Cada fila representa una interacción entre dos proteínas identificadas por sus nombres
sistemáticos (e.g., \texttt{YLR418C}, \texttt{YOL145C}). El archivo lista cada arista
en ambas direcciones, por lo que se aplica deduplicación durante la carga.

Las proteínas de \textit{Saccharomyces cerevisiae} son ampliamente usadas como
organismo modelo en biología celular. Esta red es un benchmark clásico en análisis
de redes biológicas.

\subsection{Dataset 2: Trade Network}

\begin{center}
\begin{tabular}{ll}
\toprule
\textbf{Característica} & \textbf{Valor} \\
\midrule
Tipo & Dirigido, ponderado \\
Vértices & $\approx$ 161 países \\
Aristas & Miles de relaciones comerciales \\
Peso de arista & Valor del flujo comercial (USD) \\
Formato & Pajek \texttt{.net} \\
Años disponibles & 2000, 2005, 2010, 2015, 2018 \\
\bottomrule
\end{tabular}
\end{center}

El archivo Pajek contiene una sección \texttt{*Vertices} con países y coordenadas,
seguida de \texttt{*Arcs} con las exportaciones dirigidas y su peso en USD.
Una arista $A \to B$ con peso $w$ indica que el país $A$ exportó $w$ USD al país $B$.

Para los experimentos se combinan los cinco años en un único grafo, conservando
el peso máximo observado para cada par de países. Esto produce la red de comercio
internacional más completa del dataset.
